% Bagian: computability
% Bab: computability-theory
% Subbagian: s-m-n

\documentclass[../../../include/open-logic-section]{subfiles}

\begin{document}

\olfileid[id]{cmp}{thy}{smn}
\olsection{Teorema $s$-$m$-$n$}

\begin{explain}
Teorema berikut dikenal sebagai ``teorema $s$-$m$-$n$,'' karena suatu alasan
yang sebentar lagi akan jelas. Bagian yang sulit ialah memahami apa tepatnya
yang dinyatakan teorema tersebut; setelah pernyataannya dipahami, teorema itu
akan tampak cukup jelas.
\end{explain}

\begin{thm}
\ollabel{thm:s-m-n}
  Bagi setiap pasangan bilangan asli $n$ dan~$m$, terdapat suatu fungsi
  rekursif primitif~$s^m_n$ sedemikian sehingga, bagi setiap barisan
  $e$, $a_0$, \dots, $a_{m-1}$, $y_0$, \dots, $y_{n-1}$, berlaku
  \[
  \cfind{s^m_n(e, a_0, \dots, a_{m-1})}[n](y_0, \dots, y_{n-1}) \simeq
  \cfind{e}[m+n](a_0, \dots, a_{m-1}, y_0, \dots, y_{n-1}).
  \]
\end{thm}

\begin{explain}
Akan membantu jika kita memandang $s^m_n$ sebagai fungsi yang bekerja pada
\emph{program}. Artinya, $s^m_n$ menerima suatu program~$e$ untuk fungsi
berargumen $(m+n)$, beserta masukan tetap $a_0$, \dots, $a_{m-1}$; lalu fungsi
itu mengembalikan program $s^m_n(e, a_0, \dots, a_{m-1})$ untuk fungsi
berargumen~$n$ dari argumen-argumen yang tersisa. \iftag{TMs}{Jika $e$
dipandang sebagai deskripsi suatu mesin Turing, maka
$s^m_n(e, a_0, \dots, a_{m-1})$ adalah mesin Turing yang, ketika menerima
masukan $y_0$, \dots,~$y_{n-1}$, menambahkan $a_0$, \dots,~$a_{m-1}$ di depan
untai masukan, lalu menjalankan~$e$. Dengan demikian, setiap $s^m_n$ hanyalah
fungsi rekursif primitif yang menemukan kode mesin Turing yang sesuai.}{}
\end{explain}

\end{document}
